%% fddiagram-exemplos.tex
%% Copyright 2026 Rodrigo Smarzaro
%% Parte do pacote fddiagram -- distribuido sob a LPPL 1.3c ou posterior.
%% Ver o arquivo LICENSE (ou https://www.latex-project.org/lppl.txt).
%% Repositorio e contato: https://github.com/Smarzaro/fddiagram
\documentclass[a4paper]{article}
\usepackage[utf8]{inputenc}
\usepackage[T1]{fontenc}
\usepackage[brazil]{babel}
\usepackage{geometry}
\geometry{margin=2cm}
\usepackage[color=purple]{fddiagram}

\title{Pacote \texttt{fddiagram} --- exemplos de uso}
\author{Rodrigo Smarzaro}
\date{2026-09-27}

\begin{document}
\maketitle

\section*{0. Opções padrão no carregamento do pacote}

Desde a v0.2, é possível fixar um novo padrão para \emph{todos} os
diagramas do documento já no \verb|\usepackage|. Este próprio arquivo
carrega o pacote com \verb|\usepackage[color=purple]{fddiagram}|
--- por isso todos os diagramas abaixo já nascem em roxo, mesmo sem
repetir \texttt{color=} em cada ambiente. Uma opção dada diretamente
em \verb|\begin{fddiagram}[...]| continua podendo sobrescrever esse
padrão pontualmente, só para aquele diagrama.

\section*{1. Uso básico --- universo derivado automaticamente}

Basta listar as DFs; os atributos e a ordem das colunas são deduzidos
sozinhos a partir do que aparece no texto (mesma relação da Questão 2
do gabarito de SIN220):

\begin{verbatim}
\begin{fddiagram}
  A,B -> C ;
  B,D -> E,F ;
  A,D -> G,H ;
  A -> I ;
  H -> J
\end{fddiagram}
\end{verbatim}

\begin{fddiagram}
  A,B -> C ;
  B,D -> E,F ;
  A,D -> G,H ;
  A -> I ;
  H -> J
\end{fddiagram}

\section*{2. Nomes de atributos longos}

Funciona igualmente bem com nomes de atributos por extenso (não
precisam ser uma única letra) --- aqui, a relação \textsc{Tab2} da
Questão 9, com \texttt{spacing} aumentado para caber os nomes:

\begin{verbatim}
\begin{fddiagram}[spacing=2.2]
  CodDisc -> NomeDisc, CreditosDisc ;
  CodDisc, AnoSem, SiglaTur, HoraInicio -> NumHoras
\end{fddiagram}
\end{verbatim}

\begin{fddiagram}[spacing=2.2]
  CodDisc -> NomeDisc, CreditosDisc ;
  CodDisc, AnoSem, SiglaTur, HoraInicio -> NumHoras
\end{fddiagram}

\section*{3. Universo explícito, cor e ponta de seta customizadas}

Use \texttt{attrs=} quando quiser fixar quais colunas aparecem e em
que ordem (por exemplo, para incluir um atributo que não participa de
nenhuma DF, ou para forçar uma ordem diferente da de aparição no
texto). As opções \texttt{color=} e \texttt{arrow=} controlam a
aparência:

\begin{verbatim}
\begin{fddiagram}[attrs={A,B,C,D,E}, color=blue, arrow={Triangle[open]}]
  A,B -> C ;
  C,D -> E ;
  D,E -> B
\end{fddiagram}
\end{verbatim}

\begin{fddiagram}[attrs={A,B,C,D,E}, color=blue, arrow={Triangle[open]}]
  A,B -> C ;
  C,D -> E ;
  D,E -> B
\end{fddiagram}

\section*{4. Todas as opções juntas}

\begin{verbatim}
\begin{fddiagram}[spacing=1.3, rowheight=0.8, stublen=0.25,
                  color=teal, arrow=Latex]
  F,I -> P ;
  F -> M
\end{fddiagram}
\end{verbatim}

\begin{fddiagram}[spacing=1.3, rowheight=0.8, stublen=0.25,
                  color=teal, arrow=Latex]
  F,I -> P ;
  F -> M
\end{fddiagram}

\section*{5. Rótulos alinhados à esquerda (\texttt{labelalign=left})}

Por padrão (\texttt{labelalign=local}), cada rótulo \texttt{dfN} fica
colado ao início da própria régua, formando uma "escadinha". Com
\texttt{labelalign=left}, todos os rótulos ficam alinhados numa única
coluna, na margem esquerda de todo o diagrama:

\begin{verbatim}
\begin{fddiagram}[labelalign=left]
  A,B -> C ;
  B,D -> E,F ;
  A,D -> G,H ;
  A -> I ;
  H -> J
\end{fddiagram}
\end{verbatim}

\begin{fddiagram}[labelalign=left]
  A,B -> C ;
  B,D -> E,F ;
  A,D -> G,H ;
  A -> I ;
  H -> J
\end{fddiagram}

\section*{6. Linhas-guia até o nome do atributo (\texttt{guides})}

Quando \texttt{guides=true}, cada marca (traço ou seta) ganha uma
linha fina extra, ligando-a até o nome do atributo lá em cima ---
por padrão pontilhada e cinza-claro, mas o estilo
(\texttt{guidestyle=dashed} ou \texttt{dotted}) e a cor
(\texttt{guidecolor=}) são customizáveis:

\begin{verbatim}
\begin{fddiagram}[guides=true, labelalign=left]
  A,B -> C ;
  C,D -> E ;
  D,E -> B
\end{fddiagram}
\end{verbatim}

\begin{fddiagram}[guides=true, labelalign=left]
  A,B -> C ;
  C,D -> E ;
  D,E -> B
\end{fddiagram}

\begin{verbatim}
\begin{fddiagram}[guides=true, guidestyle=dashed, guidecolor=blue!30]
  A,B -> C ;
  C,D -> E ;
  D,E -> B
\end{fddiagram}
\end{verbatim}

\begin{fddiagram}[guides=true, guidestyle=dashed, guidecolor=blue!30]
  A,B -> C ;
  C,D -> E ;
  D,E -> B
\end{fddiagram}

\section*{7. Linhas em branco e DF malformada}

Linhas em branco no meio da lista de DFs não quebram mais a
compilação, e uma DF sem (ou com mais de um) \texttt{->} agora gera
uma mensagem de erro clara, apontando a linha problemática, em vez de
travar com um erro críptico do TeX --- o restante do diagrama continua
sendo desenhado normalmente:

\begin{verbatim}
\begin{fddiagram}
  A,B -> C ;

  C -> D
\end{fddiagram}
\end{verbatim}

\begin{fddiagram}
  A,B -> C ;

  C -> D
\end{fddiagram}

\end{document}
